--- Page 1 ---
 
S1 
SUPPORTING INFORMATION 
for 
Nature of S-States in the Oxygen-Evolving Complex Resolved by High-
Energy Resolution Fluorescence Detected X-ray Absorption 
Spectroscopy 
 
Maria Chrysina, Maria Drosou, Rebeca G. Castillo, Michael Reus, Frank Neese, Vera Krewald, 
Dimitrios A. Pantazis,* Serena DeBeer* 
 
Contents: 
1. Quantification of the S-states of Samples Using EPR Spectroscopy ........................................................ 2 
2. S2 and S3 Heterogeneity ............................................................................................................................. 5 
Heterogeneity of the S2 State ..................................................................................................................... 5 
Heterogeneity of the S3 State ..................................................................................................................... 5 
3. Damage Studies ......................................................................................................................................... 7 
4. Experimental Error Evaluation .................................................................................................................. 8 
5. Quantum Chemical Cluster Model of the OEC ......................................................................................... 9 
6. Derivatives of the Pure S-state Spectra ................................................................................................... 10 
First Derivatives of the Pure S-state Spectra ........................................................................................... 10 
Second Derivatives of the Pure S-state Spectra ...................................................................................... 11 
7. Comparison with Previous Data .............................................................................................................. 12 
8. Fitting of the Pre-Edge Peaks with Voigt Curves .................................................................................... 14 
9. Si−Si−1  Differences Error ........................................................................................................................ 16 
10. Calculated Mn XAS Pre-Edge Spectra for All Computational Models ................................................ 17 
11. Assignment of the Calculated Pre-Edge XAS Spectra .......................................................................... 20 
12. Analysis of Individual Mn Ion Contributions in the S3 State Pre-edge XAS Spectra ........................... 26 
13. Example of XAS Orca 5 Input ............................................................................................................... 27 
References .................................................................................................................................................... 28 
 
 


--- Page 2 ---
 
S2 
1. Quantification of the S-states of Samples Using EPR Spectroscopy 
The progression of PS II samples to the next S-transition upon illumination is not perfect, thus 
quantification of the population of PS II centers of the sample in each S-state is essential. EPR 
spectroscopy was employed following the protocol of Messinger et al.1 The multiline signal of S2 
state at g ≈ 2 was measured in all samples in order to monitor the evolution of the S2 population 
after each visible light flash. In Figure S1, representative spectra of a series of samples after 1–5 
laser flashes are presented; they are subtraction spectra minus the S1 background. In the S1 state 
there was no sign of multiline spectrum that would be originated from any residual S2 state caused 
by the preflash. All spectra were normalized to the PS II content with respect to the gx signal of 
cytochrome c550 of the extrinsic subunit of PS II at 2300 G. The multiline intensity for each sample 
was estimated by the average of the peak-to-peak intensity of the lines indicated in Figure S1a. 
The intensity of the multiline after 1 flash, i.e. in the S2 state, was assumed to be 100% and the 
multiline after 2–5 flashes was calculated as a percentage of the intensity of S2. The average of the 
multiline intensity of all samples after each number of flashes is presented in Table 1 and in Figure 
S1b as blue circles.  
 
 
 
 
 


--- Page 3 ---
 
S3 
 
Figure S1. (a) Representative EPR spectra after 1–5 flashes. All spectra are subtracted minus the dark S1 state 
in order to eliminate the gy component from cytc550 that interferes with the multiline signal. The peaks used for 
measuring the multiline intensity after each flash are marked by dashed lines. EPR parameters: modulation 
frequency: 100 kHz, modulation amplitude: 7.5 G, microwave frequency: 9.6 GHz, microwave power: 8 mW, 
sweep time: 84 s, sweep field: 5000 G, average of 4 scans, temperature: 10 K. (b) Experimental multiline signal 
intensity after 1–5 flashes (blue circles) in comparison with the calculated multiline intensity fitted with 
parameters: miss factor = 0.06 and deficiency at S2 = 0.15. Blue circles are the average of multiline intensity of 
all samples; the standard deviation is marked by bars.  
In order to calculate the S-state content of samples after each flash a model including misses, 
double hits, initial S2 population and/or blocking at S2 was used. Different combinations of the 
above parameters were tested in order to fit the multiline intensity calculated using this model with 
the experimental; residuals of the experimental and calculated multiline intensity were used as an 
estimate of the error of the fitting. Some representative trials are presented in Table 1. First, a miss 
factor of 10% was used; higher miss factors were also tried but increased the error. It is noted that 
the same miss factor for all S-transitions was used; deviations that may exist are out of the 
resolving capability of the HERFD experiment. Consideration of double hits did not improve the 
fitting. Double hits are not expected anyway when laser is used for illumination of the samples. 
2250
2500
2750
3000
3250
3500
3750
4000
4250
4500
 
 
 Magnetic Field (G)
5 flashes
4 flashes
3 flashes
2 flashes
1 flash
0
0
1
2
3
4
5
% multiline Intensity
number of flashes
20
40
60
80
100
a
 a)
 b)


--- Page 4 ---
 
S4 
After 3 and 4 flashes there is 20% of multiline intensity that cannot be explained only by misses 
(see Table 1). Consideration of a population of PS II centers blocked at S2 state that do not progress 
further upon illumination improved the fitting. In these centers the acceptor side of PS II is 
deficient and cannot accept electrons from the OEC, to progress beyond the S2 state. A value of 
5% that is similar with that in Messinger et al. was initially used. Consideration of an initial S2 
population in dark-adapted samples led to a small improvement, but on the way smaller miss factor 
and higher population of centers blocked at S2 improved the fitting more significantly without the 
initial S2 contribution. No multiline signal could be observed in S1 spectra (not shown), consistent 
with the fact that no initial S2 population was used. Finally, the best fitting between the 
experimental and the calculated multiline intensity was achieved by using a miss factor of 6% and 
15% blockage at S2. In a separate experiment, a singly-flashed sample was illuminated with 
continuous visible light illumination at -80 °C in order to create the maximum S2 population and 
no increase of the S2 multiline signal was observed. This indicates that we have indeed ~ 100% S2 
after the first flash. The absense of delay in S-state advancement during the S1→S2 transition 
confirms that after dark adaptation the samples are synchronized in the S1 state with 100% oxidized 
YD. The blockage at S2 state is too high. Exogenous electron acceptor has been added but it may 
not be integrated well in all PS II centers. To the end, the population of each S-state after each laser 
flash is presented in Table 2. After “0, 1, 2, 3 flashes” there is 100% S1, 94% S2, 74% S3, and 70% 
S0 respectively.  
Table S1. Calculated multiline intensity in comparison with the experimental intensity of the multiline signal of 
the S2 state using different combinations of parameters. All parameters are given as percentages (%). 
 
misses 
10 
10 
10 
10 
10 
7 
6 
 
 
 
double hits 
0 
5 
0 
0 
0 
0 
0 
 
 
 
impair at S2 
0 
0 
5 
5 
10 
10 
15 
 
 
 
initial S2 
0 
0 
0 
5 
5 
5 
0 
 
 
calculated multiline 
intensity 
1 flash 
100 
100 
100 
100 
100 
100 
100 
experimental 
multiline intensity 
100 
2 flashes 
20 
20 
25 
25 
30 
24 
27 
26 
3 flashes 
3 
4 
9 
8 
14 
12 
17 
19 
4 flashes 
0 
15 
6 
10 
15 
15 
16 
14 
5 flashes 
66 
60 
68 
69 
71 
79 
82 
82 
 
Error1 
744 
771 
384 
296 
163 
63 
9 
 
 
1 The error of each fitting was calculated as the sum of residuals between calculated and experimental multiline intensity, as in Messinger et al.1 


--- Page 5 ---
 
S5 
Table S2. Calculated S-state content after a given number of flashes. 
# flashes 
S1 
S2 
S3 
S0 
0 
100 
 
 
 
1 
6 
94 
 
 
2 
1 
25 
74 
 
3 
0 
16 
14 
70 
4 
65 
15 
2 
18 
5 
20 
77 
0 
3 
 
 
2. S2 and S3 Heterogeneity 
Heterogeneity of the S2 State 
It was mentioned in the main text that we avoided using glycerol in the samples in order to achieve 
a homogeneous S3 that has been already characterized by EPR, hyperfine spectroscopy and DFT.2 
This choice has the drawback that a minority of centers of S2 state represents the high spin form. 
In order to quantify this minor S2 population of our samples we compared the present S2 with 
methanol treated S2 (Figure 2). Methanol converts the minor component of S2 (high spin) to the 
major component (low spin) that is represented by the multiline signal. In the methanol treated 
sample, an increase in the intensity of the S2 multiline signal is observed by 20%, which implies 
that the untreated sample contains 20% of the minor high-spin component of S2. Thus, the S2 state 
of this study represents the low-spin form in 80% of centers. 
 
Heterogeneity of the S3 State 
W-band EPR/EDNMR and DFT revealed an open cubane structure of the S3 state with four 6-
coordinate MnIV in the thermophilic cyanobacterium T. vestitus; a new water that is not present in 
S2 is bound to Mn1.2 The EPR spectrum is shown in green in Figure S2b. Three perturbed 
configurations of this motif were observed in glycerol treated PS II from cyanobacterium 
Synechocystis sp by 130 GHz EPR.3 A form of the S3 state with a Mn ion to represent highly 


--- Page 6 ---
 
S6 
anisotropic hyperfine couplings was observed by W-band EPR/EDNMR in methanol or glycerol 
treated or Sr-substituted samples from T. vestitus.4 The aforementioned S3 forms represent a total 
spin of S = 3. A high spin S3 form (S = 6) was observed by Q-band EPR spectroscopy in methanol 
treated PS II from spinach and also as a major component in untreated PS II,5 as first predicted by 
quantum chemical studies.6  
In Figure S2b representative spectra of the two configurations of the S3 state in T. vestitus are 
presented. The first observed high-field EPR signal of the S3 (“narrow”) is shown in green. Upon 
addition of glycerol or methanol in the samples, a “wide” component is observed (marked with the 
red arrow) together with the “narrow” signal. The populations of each component depend on the 
percentage of methanol of glycerol added. We used glycerol-free samples that represent only the 
“narrow” S3 (green in Figure S2b).  
 
Figure S2. Inhomogeneity of the S2 (a) and S3 (b) states as indicated by X-band and W-band EPR spectroscopy 
respectively. The green spectra in both figures are the homogeneous S2 and S3 configurations, while the black 
represent mixtures (see text). In the samples used herein the S2 is inhomogeneous (black in panel a), while the 
S3 is homogeneous (green in panel b). (a) Comparison of the S2 state in untreated and methanol-treated PS II. 
The 20% increase of the multiline intensity upon addition of methanol indicates that 20% of the population in 
the untreated S2 represents a high spin form, not observable by X-band. EPR parameters as in Figure S1. (b) W-
band ESE-detected field-swept spectra of the untreated S3 (green), S3 with 10% glycerol and 3% methanol 
(black). Cytochrome signals are marked with grey lines. The untreated S3 signal (narrow) is marked by a green 
arrow while the recently revealed configuration of S3 (wide) by a red arrow. Panel S2b adapted from Chrysina 
et al. PNAS 2019, 116 (34), 16841-16846. Copyright 2019 National Academy of Sciences 
 
2250 2500 2750 3000 3250 3500 3750 4000 4250 4500
 
 
 
Magnetic Field (G)
untreated
4% v/v
methanol
a) S2
1.5          2.0          2.5          3.0           3.5          4.0          4.5         5.0
Magnetic Field (T)
10% v/v 
glycerol
3% v/v
methanol
b) S3
cytc550
untreated
cytc550


--- Page 7 ---
 
S7 
3. Damage Studies 
Before starting the actual measurement of each sample (no flash, 1, 2, 3 flashes), damage studies 
were performed (described in methods section). Here, some representative spectra are presented. 
 
Figure S3. Damage test in 0, 1, 2, 3 flashed samples: 2 rounds of measurement were performed on the same 
spots. In all cases flux was attenuated to 6%. For 0 and 1 flashed sample (S1 and S2 samples): scan time was 58 
s, average of three different spots was used. No damage was observed in 116 s. In 2 and 3 flashed samples, each 
scan is 38 s. No damage was observed in 76 s. For the 2-flashed sample (mainly S3 state sample) the average of 
37 spots was tested for damage and for the 3-flashed sample (mainly S0 state sample), the average of 8 spots.  
 
 
6535 6540
6545
6550
6555
6560
6565
6570
 
 
 
 Energy (eV)
no flash
 
1flash
 
 
 
2 flashes
 
 
 
3 flashes
 
 
6535 6540
6545
6550
6555
6560
6565
6570
6535 6540
6545
6550
6555
6560
6565
6570
6535 6540
6545
6550
6555
6560
6565
6570
 Energy (eV)
 Energy (eV)
 Energy (eV)
1
st
 round
2
nd
 round
1
st
 round
2
nd
 round
1
st
 round
2
nd
 round
1
st
 round
2
nd round


--- Page 8 ---
 
S8 
4. Experimental Error Evaluation 
 
Figure S4. Standard error of the 0, 1, 2, 3 flashed spectra (left) and of the pure S1, S2, S3, S0 spectra (right) is 
represented as a shaded area. Standard error was calculated as described in Material and Methods.  
 
 
 
 
 
 
6535
6540
6545
6550
6555
6560
6565
6570
0.0
0.2
0.4
0.6
0.8
1.0
1.2
Normalised Intensity
Energy (eV)
S1
S2
S3
S0
 
 
6537 6538 6539 6540 6541 6542 6543 6544 6545 6546 6547
0.00
0.02
0.04
0.06
0.08
0.10
0.12
Normalised Intensity
Energy (eV)
 
 
0 flash
1 flash
2 flashes
3 flashes
6535
6540
6545
6550
6555
6560
6565
6570
0.0
0.2
0.4
0.6
0.8
1.0
1.2
Normalised Intensity
Energy (eV)
6537 6538 6539 6540 6541 6542 6543 6544 6545 6546 6547
0.00
0.02
0.04
0.06
0.08
0.10
0.12
Normalised Intensity
Energy (eV)
0 flash
1 flash
2 flashes
3 flashes
S1
S2
S3
S0


--- Page 9 ---
 
S9 
5. Quantum Chemical Cluster Model of the OEC 
 
 
Figure S5. QM Model of the S1 state. H atoms attached to carbons are omitted for clarity. 
 
 


--- Page 10 ---
 
S10 
6. Derivatives of the Pure S-state Spectra 
First Derivatives of the Pure S-state Spectra 
Figure S6. First derivative (with no smoothing at left and 41 point smoothing at right) of the XANES spectra (5 
point smoothed). The XANES spectra were 5-point smoothed, the first derivative was calculated and the first 
derivative was 41-point smoothed. 
 
 
6535 6540 6545 6550 6555 6560 6565 6570
-0.1
-0.05
0
0.05
0.1
0.15
0.2
0.25
0.3
S0
 
 
   
6535 6540 6545 6550 6555 6560 6565 6570
-0.1
-0.05
0
0.05
0.1
0.15
0.2
0.25
0.3
S1
S2
S3
S0
 Energy (eV)
 Normalised Intensity
 Energy (eV)


--- Page 11 ---
 
S11 
Second Derivatives of the Pure S-state Spectra 
 
Figure S7. Second derivative of the HERFD data (left) in comparison to Messinger et al.1 second derivative 
(right). 
 
 


--- Page 12 ---
 
S12 
7. Comparison with Previous Data 
 
Figure S8. HERFD spectra of this work (black) in comparison with data reported from Messinger et al.1 (red). 
Adapted in part from Messinger at al. J. Am. Chem. Soc. 2001, 123, 7804-7820. Copyright 2001 American 
Chemical Society.  
 
 
6535
6540
6545
6550
6555
6560
6565
6570
0.0
0.2
0.4
0.6
0.8
1.0
1.2
6535
6540
6545
6550
6555
6560
6565
6570
0.0
0.2
0.4
0.6
0.8
1.0
1.2
6535
6540
6545
6550
6555
6560
6565
6570
0.0
0.2
0.4
0.6
0.8
1.0
1.2
6535
6540
6545
6550
6555
6560
6565
6570
0.0
0.2
0.4
0.6
0.8
1.0
1.2
6538 6540 6542 6544 6546 6548
6538 6540 6542 6544 6546 6548
6538 6540 6542 6544 6546 6548
6538 6540 6542 6544 6546 6548
 
 
Normalised Intensity
S0
S3
S1
 
 
 Normalised Intensity
 
 
Energy (eV)
S2
 
 
 
Energy (eV)


--- Page 13 ---
 
S13 
 
Figure S9. Left: Comparison between XAS spectra of Haumann et al.7 (in red) and Messinger et al.1 (in black) 
Right: Comparison between XAS spectra of Haumann et al.7 (in red) and HERFD spectra of this work (in black). 
Adapted in part from Messinger at al. J. Am. Chem. Soc. 2001, 123, 7804-7820 (Copyright 2001 American 
Chemical Society) and Haumann et al Biochemistry 2005, 44, 1894-1908 (Copyright 2005 American Chemical 
Society). 
Haumann et al.7 data are in spinach, 1M betaine as a cryoprotectant and 10%v/v glycerol to 
maximize the multiline S2 signal. No g = 4.1 was reported to be observed, thus S2 is assumed to 
be 100% in the low-spin form. Partially dehydrated samples. Messinger et al.1 experiments are in 
spinach, 30% glycerol. There are no high field EPR data in spinach S3 with glycerol at the moment. 
6535 6540 6545 6550 6555 6560 6565 6570
0
1
S0
S1
S2
 Energy (eV)
 
S3
Haumann et al.
Messinger et al.
 Normalised Intensity
0
1
 
 
6535 6540 6545 6550 6555 6560 6565 6570
 Energy (eV)
Haumann et al.
HERFD (this work)


--- Page 14 ---
 
S14 
 
8. Fitting of the Pre-Edge Peaks with Voigt Curves 
 
Figure S10. Fitting of the pre-edge experimental peaks (black) with Voigt curves (red). The individual curves 
used for the fitting are presented in blue and the residuals in orange. Fitted parameters are listed in Table S3. 
 
 
S1
S0
6538 6539 6540 6541 6542 6543 6544 6545 6546 6547 6548
-0.01
0
0.01
0.02
0.03
0.04
0.05
0.06
0.07
0.08
0.09
Energy (eV)
Normalized Intensity (a.u.)
S1
0.10
S1
6538 6539 6540 6541 6542 6543 6544 6545 6546 6547 6548
-0.01
0
0.01
0.02
0.03
0.04
0.05
0.06
0.07
0.08
0.09
0.10
6538 6539 6540 6541 6542 6543 6544 6545 6546 6547 6548
-0.01
0
0.01
0.02
0.03
0.04
0.05
0.06
0.07
0.08
0.09
0.10
S0
6538 6539 6540 6541 6542 6543 6544 6545 6546 6547 6548
-0.01
0
0.01
0.02
0.03
0.04
0.05
0.06
0.07
0.08
0.09
0.10
Normalized Intensity (a.u.)
Normalized Intensity (a.u.)
S2
S3
Energy (eV)
6538 6539 6540 6541 6542 6543 6544 6545 6546 6547 6548
-0.01
0
0.01
0.02
0.03
0.04
0.05
0.06
0.07
0.08
0.09
0.10
6538 6539 6540 6541 6542 6543 6544 6545 6546 6547 6548
-0.01
0
0.01
0.02
0.03
0.04
0.05
0.06
0.07
0.08
0.09
0.10
exp
fitting sum
fitting
residuals


--- Page 15 ---
 
S15 
 
 
Table S3. Parameters of the fitting of the pre-edge region. 
 
Amplitude      Energy        Gauss        Lorentz 
                         (eV)          FWHM      FWHM 
Amplitude      Energy        Gauss        Lorentz 
                         (eV)          FWHM      FWHM 
S0 
Fitting 1 (error = 280.2 10-6) 
Fitting 2 (error = 455.4 10-6) 
peak 1 
peak 2 
peak 3 
peak 4 
0.030          6539.864       1.007          0.000 
0.082          6540.866       1.549          0.000 
0.080          6542.455       1.549          0.100 
0.017          6545.336       1.085          0.084 
   0.100         6540.461       1.700          0.000 
   0.092         6542.326       1.700          0.100 
   0.017         6545.338       1.097          0.100 
S1 
Fitting 1 (error = 324.5 10-6) 
Fitting 2 (error = 190.2 10-6) 
peak 1           0.119          6540.799       1.700          0.027 
peak 2           0.124          6542.651       1.683          0.463 
   0.008     6539.681       0.806          0.000 
   0.118     6540.896       1.692          0.000 
   0.115     6542.703       1.607          0.435 
S2 (error = 226.8 10-6) 
 
peak 1           0.108         6540.816       1.527          0.001 
 
peak 2           0.108         6542.802       1.700          0.000 
 
S3 (error = 449.1 10-6) 
 
peak 1            0.113        6540.994       1.670          0.000 
 
peak 2            0.122        6542.878       1.573          0.100 
 
 
Two options are presented for the fitting of S0 and S1 states with Voigt curves. For S0, the error is 
lower for fitting 3 peaks in the range 6538-6545 eV rather than 2 peaks (intensity of residuals 
lower). In the case of S1, the error is lower for fitting 3 peaks in this range; the one small intensity 
peak fits a shoulder at 6539.5 eV. The error of the fitting for the S3 state is high but mainly because 
of the range > 6545 eV (higher intensity of residuals). 
 
 


--- Page 16 ---
 
S16 
9. Si−Si−1  Differences Error 
 
Figure S11. Standard error of the Si−Si−1 differences.  
 
 
-0.05
0.00
0.05
 0.05
0.00
6538 6539 6540 6541 6542 6543 6544 6545
-0.05
0.00
 
 
 
S1 - S0
 
 
 
Normalised Intensity
S2 - S1
 
 
Energy (eV)
S3 - S2


--- Page 17 ---
 
S17 
10. Calculated Mn XAS Pre-Edge Spectra for All Computational Models 
 
 
Figure S12. Schematic depiction of the inorganic cores of all S0–S3 models examined in this work, with the Mn 
oxidation states indicated in purple for Mn(IV) and pink for Mn(III), and the Jahn–Teller elongation axis of each 
Mn(III) ion shown with thicker pink lines. Selected optimized distances (in Å, Mn1/Mn4–O5 in black and O5–
O6/W2 in red) are indicated to facilitate comparisons. Complete coordinates of all models are provided in full 
as Supporting Information.  
 


--- Page 18 ---
 
S18 
 
Figure S13. Calculated XAS spectra of all models for each state shown in different colors compared to 
experiment (black curves). A 1.1 eV broadening and a constant shift of 36.3 eV have been applied to the 
calculated spectra. 
 
 


--- Page 19 ---
 
S19 
 
Figure S14. Comparison of experimental S2−S1 data (black line) with the calculated S2−S1 difference spectra 
for models 𝐒𝟏
 𝐁,𝐇 and 𝐒𝟐
 𝐀,𝐇 (blue line) and with the calculated S2−S1 difference spectra for model 𝐒𝟏
 𝐁,𝐇 and a 
mixture of 80% 𝐒𝟐
 𝐀,𝐇 and 20% 𝐒𝟐
 𝐁,𝐇 (red line). The effect of 20% 𝐒𝟐
 𝐁,𝐇 admixture on the difference spectra is 
negligible and both S2−S1 difference spectra reproduce experiment.  
 
 
Figure S15. Experimental and calculated S1–S0 spectral differences for different of S0 and S1 isomers.  
 
 
 


--- Page 20 ---
 
S20 
11. Assignment of the Calculated Pre-Edge XAS Spectra 
 
 
Figure S16. Assignment of the calculated pre-edge XAS spectrum based on the NTOs associated with the 
transitions for models 𝐒𝟏
 𝐁 and 𝐒𝟏
 𝐁,𝐇. Assignment of transitions for 𝐒𝟏
 𝐁,𝐇 is similar to that of 𝐒𝟏
 𝐁. The local Mn4 
1s → 3d transition at ~6540 eV is more intense in 𝐒𝟏
 𝐁,𝐇. 
 


--- Page 21 ---
 
S21 
 
Figure S17. Assignment of the calculated pre-edge XAS spectrum based on the NTOs associated with the 
transitions for model 𝐒𝟐
 𝐀. 
 


--- Page 22 ---
 
S22 
 
Figure S18. Assignment of the calculated pre-edge XAS spectrum based on the NTOs associated with the 
transitions for model 𝐒𝟑
 𝐀,𝐖. 
 


--- Page 23 ---
 
S23 
 
Figure S19. Assignment of the calculated pre-edge XAS spectrum based on the NTOs associated with the 
transitions for model 𝐒𝟑
 𝐨𝐱𝐲𝐥.𝐨𝐱𝐨. 
 


--- Page 24 ---
 
S24 
 
Figure S20. Assignment of the calculated pre-edge XAS spectrum based on the NTOs associated with the 
transitions for model 𝐒𝟑
 𝐩𝐞𝐫𝐨𝐱𝐨. 
 


--- Page 25 ---
 
S25 
 
Figure S21. Assignment of the calculated pre-edge XAS spectrum based on the NTOs associated with the 
transitions for model 𝐒𝟎
 𝐎𝟒. 
 
 
 
 
 
 
 
 
 


--- Page 26 ---
 
S26 
12. Analysis of Individual Mn Ion Contributions in the S3 State Pre-edge XAS 
Spectra 
 
Figure S22. Calculated XAS spectra for individual Mn ions for 𝐒𝟑
 𝐀,𝐖, 𝐒𝟑
 𝐁,𝐖, 𝐒𝟑
 𝐨𝐱𝐲𝐥.𝐨𝐱𝐨 and 𝐒𝟑
 𝐩𝐞𝐫𝐨𝐱𝐨. The higher 
intensity of the transitions of the calculated 𝐒𝟑
 𝐨𝐱𝐲𝐥.𝐨𝐱𝐨 spectra compared to the 𝐒𝟑
 𝐀,𝐖 in the 6541–6544 eV region, 
is attributed mostly to excitations from Mn1 and Mn4 1s core orbitals, while the calculated 𝐒𝟑
 𝐩𝐞𝐫𝐨𝐱𝐨 spectra show 
lower overall intensity than 𝐒𝟑
 𝐀,𝐖 mostly due to excitations from Mn1 and Mn3 1s core orbitals. 
 
 


--- Page 27 ---
 
S27 
13. Example of XAS Orca 5 Input  
 
! UKS TPSSh RIJCOSX ZORA ZORA-def2-TZVP(-f) SARC/J CPCM 
! NoTrah TightSCF 
 
%pal nprocs 10 end 
%maxcore 15000 
%basis newgto C "ZORA-def2-SVP" end 
       newgto H "ZORA-def2-SVP" end 
end 
%cpcm 
    surfacetype vdw_gaussian 
    epsilon 6.0 
end 
%scf maxiter 200 
     shift shift 0.10 erroff 0.1 end 
     FlipSpin 1,2 FinalMs 0.5 
end 
%tddft NRoots 150 
        MaxDim 8 
        OrbWin[0]=0,0,-1,-1 
        OrbWin[1]=0,0,-1,-1 
        DoQuad true 
End 
 
*xyzfile 1 14 filename.xyz 
 
 
 


--- Page 28 ---
 
S28 
References 
(1) 
Messinger, J.; Robblee, J. H.; Bergmann, U.; Fernandez, C.; Glatzel, P.; Visser, H.; Cinco, R. M.; 
McFarlane, K. L.; Bellacchio, E.; Pizarro, S. A.; Cramer, S. P.; Sauer, K.; Klein, M. P.; Yachandra, V. K. 
Absence of Mn-centered oxidation in the S2 → S3 transition: implications for the mechanism of 
photosynthetic water oxidation. J. Am. Chem. Soc. 2001, 123, 7804-7820. 
(2) 
Cox, N.; Retegan, M.; Neese, F.; Pantazis, D. A.; Boussac, A.; Lubitz, W. Electronic structure of 
the oxygen-evolving complex in photosystem II prior to O-O bond formation. Science 2014, 345, 804-808. 
(3) 
Marchiori, D. A.; Debus, R. J.; Britt, R. D. Pulse EPR Spectroscopic Characterization of the S3 
State of the Oxygen-Evolving Complex of Photosystem II Isolated from Synechocystis. Biochemistry 2020, 
59, 4864–4872. 
(4) 
Chrysina, M.; Heyno, E.; Kutin, Y.; Reus, M.; Nilsson, H.; Nowaczyk, M. M.; DeBeer, S.; Neese, 
F.; Messinger, J.; Lubitz, W.; Cox, N. Five-coordinate MnIV intermediate in the activation of nature’s water 
splitting cofactor. Proc. Natl. Acad. Sci. U. S. A. 2019, 116, 16841-16846. 
(5) 
Zahariou, G.; Ioannidis, N.; Sanakis, Y.; Pantazis, D. A. Arrested Substrate Binding Resolves 
Catalytic Intermediates in Higher-Plant Water Oxidation. Angew. Chem., Int. Ed. 2021, 60, 3156-3162. 
(6) 
Retegan, M.; Krewald, V.; Mamedov, F.; Neese, F.; Lubitz, W.; Cox, N.; Pantazis, D. A. A five-
coordinate Mn(IV) intermediate in biological water oxidation: spectroscopic signature and a pivot 
mechanism for water binding. Chem. Sci. 2016, 7, 72-84. 
(7) 
Haumann, M.; Muller, C.; Liebisch, P.; Iuzzolino, L.; Dittmer, J.; Grabolle, M.; Neisius, T.; Meyer-
Klaucke, W.; Dau, H. Structural and oxidation state changes of the photosystem II manganese complex in 
four transitions of the water oxidation cycle (S0 → S1, S1 → S2, S2 → S3, and S3,4 → S0) characterized by 
X-ray absorption spectroscopy at 20 K and room temperature. Biochemistry 2005, 44, 1894-1908.